% problem-000146 5 分子机器氧化还原
\section*{5. 分子机器氧化还原}
（图：）

\subsection*{5-1}
指出DB和CBPQT分⼦所带电荷数（表明正负）。

\subsection*{5-2}
通常条件下DB是⽆法穿⼊到CBPQT的环中，CBPQT在经过2个电⼦还原形成⾃由基之后，可以被DB的左端穿⼊，与经过1个电⼦还原之后的⽚段V结合，简述其原因并画出CBPQT在接受了2个电⼦之后的还原产物最稳定的结构简式并注明电荷和⾃由基的位置。

\subsection*{5-3}
将上述结构中的CBPOT与⽚段V都重新氧化，CPBQT会进⼀步翻越⽚段IPP，像“泵”⼀样被推⾄⽚段T和⽚段S之间的区域，且不会脱落，形成下图中的产物。指出⽚段IPP上的异丙基所起的作⽤。

（图：）
PY

（图：）

（图：）

（图：）
T
图 5-2

（图：）

（图：）
S

\subsection*{5-4}
分⼦泵通过⼀组“还原-氧化”过程，推动了标准吉布斯⾃由能变远⼤于0的化学反应的发⽣，起到了对化学反应“做功”的作⽤。整个DB穿⼊CBPQT环中的反应的标准吉布斯⾃由能变为59 kJ $\mathrm { m o l ^ { - 1 } }$ 。反应过程表⽰如下：

$
\mathrm{CBPQT} + \mathrm{DB} \rightarrow \mathrm{CBPQT-DB}
$

反应第⼀步采⽤ Zn 作为还原剂（标准电极电势为−0.76 V），将 CBPQT 和 DB 还原为[CBPQT-DB]3−。反应第⼆步采⽤ MB+作为氧化剂（标准电极电势为 1.22 V），将[CBPQT-DB]3−氧化为 CBPQT-DB，MB+则转变为 MB。DB/DB−和 CBPOT/CBPOT2−的标准电极电势分别为−0.32 V 和−0.26 V，计算分⼦泵的理论⼯作效率�。
